\documentclass[11pt,a4paper]{article}



\usepackage[T2A]{fontenc}
\usepackage[utf8]{inputenc}
\usepackage[english,russian]{babel}


\begin{document}

%\input{Sections/top}
\section{План курса}
Курс посвящён современным методам управления, особенно актуальным для робототехники, промышленности, транспорта и приложений ИИ. Мы охватываем темы от сертификатов устойчивости до SOS, MPPI и автономного управления роботами.

\section{Система оценивания}
Критерии оценок:
\begin{itemize}
	\item A: 85–100 баллов
	\item B: 70–84 балла
	\item C: 50–69 баллов
	\item D: менее 50 баллов
\end{itemize}

Распределение баллов:
\begin{itemize}
	\item 20 баллов можно получить за Задание 1
	\item 20 баллов можно получить за Задание 2
	\item 10 баллов можно получить за активность на занятиях
	\item 50 баллов выставляются за экзамен
\end{itemize}

\subsection{Политика опозданий}
Баллы могут быть снижены, если работа сдана после крайнего срока.

\section{Структура курса}

Мы изучаем следующие темы:

\begin{enumerate}
	\item Введение, повторение теории управления
	\item Сертификаты нелинейной устойчивости, область притяжения
	\item Купман, SOS (для сертификатов устойчивости)
	\item Базовое нелинейное управление (компенсация гравитации, CTC, адаптивное управление в регрессорной форме)
	\item Нелинейное MPC
	\item MPPI
	\item Продвинутое линейное управление: задержка времени, наблюдатель возмущений
	\item Продвинутое линейное управление: робастное MPC
	\item Расширенный фильтр Калмана
	\item Планирование траекторий: коллокация, стрельба
	\item Управление на основе обучения с подкреплением
	\item Управление шагающими роботами: модель Райберта, MPC Кима, WBIC
	\item Управление шагающими роботами: наблюдатель на основе групп Ли
	\item Управление БПЛА: дифференциальная плоскность
\end{enumerate}

\section{План курса}

\subsection{Области знаний (в терминах приложений)}

\begin{itemize}
	\item Робототехника
	\item Автоматизация
	\item Проектирование (механическое, электрическое)
	\item Управление
	\item Анализ данных
	\item Вычислительная геометрия
\end{itemize}

\subsection{Формат проведения}
Лекции проводятся каждую неделю, за ними следуют практические занятия (лабораторные работы). Во время занятий студенты должны разрабатывать собственный код на основе знаний, полученных на лекциях и при самостоятельном изучении.

\subsection{Предварительные требования}

\begin{itemize}
	\item Строгие требования: Линейная алгебра.
	\item Слабые требования: Математический анализ, Теория управления.
	\item Необходимые фоновые знания: Python (допустимы Matlab или любой другой язык, подходящий для работы с задачами, требующими линейной алгебры).
\end{itemize}

\subsection{Ожидаемые результаты обучения}

Курс предоставит участникам возможность:
\begin{itemize}
	\item Понять и научиться использовать современные методы управления: прямой метод Ляпунова, сертификаты устойчивости на основе сумм квадратов, представление Купмана и др.
	\item Научиться применять современные методы управления в робототехнике.
	\item Уметь реализовывать различные алгоритмы, рассматриваемые в курсе (см. темы).
\end{itemize}

\subsection{Рекомендуемые материалы}
\begin{itemize}
	\item Аннотированные слайды
	\item Онлайн-материалы
	\item Учебные видео
\end{itemize}

\subsection{Компьютерные ресурсы}
Студентам потребуется проводить компьютерные эксперименты на ноутбуке и/или на лабораторных компьютерах.

\subsection{Лабораторные упражнения}
Для курса подготовлена серия лабораторных работ и электронных раздаточных материалов.

\subsection{Лабораторные ресурсы}
Студентам потребуется использовать и модифицировать программный инструмент, написанный на Python, который работает на нескольких платформах (Linux, Microsoft Windows и Mac OS). Инструмент требует свободно распространяемых библиотек.

\subsection{Политика сотрудничества и цитирования}
Мы поощряем интенсивное обсуждение и сотрудничество в этом классе. Вы можете свободно обсуждать все аспекты курса с одногруппниками и работать вместе над выполнением заданий и отчёта по проекту. Однако, если вы работаете совместно, вы должны указать детали вашего вклада и вклада других участников.
\end{document}












